\documentclass[11pt]{article}
\usepackage[margin=1in]{geometry}
\usepackage{amsmath}
\usepackage{booktabs}
\usepackage{siunitx}
\usepackage{enumitem}
\usepackage[colorlinks=true, linkcolor=blue, urlcolor=blue]{hyperref}
\usepackage{xcolor}

\newcommand{\code}[1]{\texttt{#1}}
\newcommand{\gib}{\,\text{GiB}}

\title{The LLMHub GPU Fit Estimator:\\ Memory Model, Launch Gate, and Production Hardening}
\author{LLMHub / SPIN}
\date{August 14, 2026}

\begin{document}
\maketitle

\begin{abstract}
LLMHub now estimates whether a vLLM deployment will fit on NCSA Delta GPUs
before submitting it to Slurm, and blocks launches that provably cannot boot.
This note explains the memory model and its calibration, the launch-gate
contract, the five production defects found during pre-release review (two
false-accept channels, three availability channels), the fixes, and the
remaining caveats the team should know about.
\end{abstract}

\section{Two verdicts, one model}

Everything in the estimator rests on distinguishing two questions that sound
alike but have very different answers:

\begin{description}[leftmargin=1.6em]
\item[``Does it fit?'' (\code{fits})] Could the GPU hold the \emph{worst
  case} --- weights, plus KV cache for \code{max\_num\_seqs} concurrent
  sequences each at full context, plus overhead? This is a planning number
  shown in the UI survey. vLLM never actually reserves this much.
\item[``Does it boot?'' (\code{starts})] After loading weights and overhead,
  is the leftover KV pool big enough to hold \emph{one} full-length sequence?
  This is the check vLLM itself performs at startup, and it is what the launch
  gate enforces. Beyond the pool, vLLM queues and preempts requests --- it does
  not OOM --- so sustained concurrency is reported as capacity, never gated.
\end{description}

\noindent
\code{fits}${}=$\code{false} with \code{starts}${}=$\code{true} is the
\emph{normal} operating regime, not an edge case: Qwen2.5-7B on one A40 at
context 4096 and concurrency 256 ``needs'' a fictional 76\gib{} but boots with
a real 24.9\gib{} KV pool ($\approx$466k tokens, $\approx$113 concurrent
full-context sequences).

\section{The memory model}

All quantities are per GPU, in binary GiB ($2^{30}$ bytes). Three additive
terms are compared against measured VRAM:
\begin{equation}
\underbrace{\frac{W}{t}}_{\text{weights}}
\;+\;
\underbrace{n_{\text{tok}} \cdot k_{\text{tok}}(t)}_{\text{KV cache}}
\;+\;
\underbrace{(1-u)\,V \;+\; 0.8 \;+\; 0.002\,m \;+\; 0.25\,[t>1]}_{\text{overhead}}
\;\le\; V
\end{equation}
where $W$ is the model's weight size (read from safetensors metadata, not
computed), $t$ the tensor-parallel size, $u=0.9$ the vLLM memory-utilization
fraction, $V$ the measured per-GPU VRAM, $m$ the scheduler concurrency
(\code{--max-num-seqs}), and the per-token KV cost is
\begin{equation}
k_{\text{tok}}(t) = 2 \cdot L \cdot \max\!\Big(1,\ \frac{H_{kv}}{t}\Big) \cdot d \cdot b,
\end{equation}
with $L$ layers, $H_{kv}$ KV heads, $d$ head dimension, and $b$ bytes per
element. The $\max(1,\cdot)$ term captures vLLM's KV-head \emph{replication}:
when $H_{kv} < t$ (common for GQA at high TP), per-GPU KV stops shrinking.
Weights shard by $t$; the utilization reserve and internal overhead are paid in
full by \emph{every} rank.

\subsection{Calibration}

The internal-overhead line $0.8 + 0.002\,m$ was fitted as a conservative upper
bound over eight real vLLM~0.11.0 startup probes on Delta (Qwen2.5-0.5B/7B;
A40 and A100-40GB; TP 1/2/4; $m \in \{32,256,512,1024\}$), reading vLLM's own
``Model loading took'' and ``Available KV cache memory'' log lines against
\code{nvidia-smi} totals. The model under-promises the measured KV pool at
every probe point:

\begin{center}
\begin{tabular}{lcccc}
\toprule
Probe & $m$ & Predicted pool & Measured & Margin \\
\midrule
0.5B / A40 / TP1 & 256 & 38.25 & 38.51 & $+0.26$ \\
7B / A40 / TP1 & 256 & 24.93 & 25.17 & $+0.24$ \\
7B / A40 / TP2 & 256 & 31.81 & 32.26 & $+0.45$ \\
7B / A40 / TP4 & 256 & 35.37 & 35.85 & $+0.48$ \\
7B / A100 / TP1 & 256 & 20.44 & 20.73 & $+0.29$ \\
7B / A40 / TP1 & 32 & 25.38 & 25.43 & $+0.05$ \\
7B / A40 / TP1 & 512 & 24.42 & 24.66 & $+0.24$ \\
7B / A40 / TP1 & 1024 & 23.39 & 23.64 & $+0.25$ \\
\bottomrule
\end{tabular}
\end{center}

\noindent
A regression test asserts \code{predicted}~$\le$~\code{measured} with zero
slack (safety: an optimistic model would let the gate approve launches that
OOM) plus an accuracy bound of 1.2\gib. VRAM values in the hardware table are
\code{nvidia-smi} measurements, not marketing numbers (A40: 44.988, not 48).

\section{The launch gate}

Immediately before vec-inf submits a Slurm job, the gate resolves the launch
tuple --- HF model id, partition, context, TP, concurrency --- with the
precedence that actually reaches vLLM: user free-form \code{vllm\_args}
$>$ explicit request fields $>$ curated catalog (tensor-parallel size is the
exception --- the launch path never emits it from the request's GPU count, so
the catalog flag wins there). It then certifies the boot contract:
\begin{equation}
\frac{W}{t} + C \cdot k_{\text{tok}}(t) + \text{overhead}(m_{\text{launch}}) \;\le\; V ,
\end{equation}
where $C$ is the resolved \code{max\_model\_len} and, crucially,
\emph{the KV budget is one sequence but the overhead is charged at the
concurrency the job will actually boot with}
($m_{\text{launch}} = $ user override $>$ catalog $>$ the vLLM default ---
\textbf{1024} on the V1 engine in the production container, not the V0-era
256; see \S\ref{sec:live}).

The gate \textbf{blocks} only on a real sizing verdict: the pool cannot hold
one full-context sequence, or TP exceeds the partition's GPU count. It
\textbf{skips with a warning} --- launch proceeds, auditable in logs --- whenever
it cannot give an honest verdict: partition not in the Delta hardware table,
multi-node launches (not modeled), non-NVIDIA GPUs, unresolvable model
metadata (sparse VLM configs, HF outage, gated repo), or a failed catalog
lookup. The public \code{POST /api/validate-config} endpoint remains strictly
fail-closed (``cannot verify'' is a refusal there); the skip semantics apply
only to the launch gate, where blocking a curated model that launched fine
yesterday would be the worse failure. An estimator crash anywhere in this path
is caught and logged; it never blocks a launch or leaks the GPU allocation.

\section{What pre-release review found and fixed}

\subsection{False-accept: overhead charged at the wrong concurrency}

The gate originally charged overhead at $m=1$ while 90 of 94 catalog entries
boot at the vLLM default $m=256$. The $0.002 \cdot 255 \approx 0.51\gib$ gap
exceeds every calibration margin (0.24--0.48\gib{} at $m=256$), so a genuine
false-accept window existed. Concretely, for Llama-3.3-70B on
\code{gpuA100x4} (TP=4): the $m{=}1$ contract certifies a pool of
$\approx$27{,}400 tokens, while the calibrated pool at $m{=}256$ holds
$\approx$20{,}800 --- a launch at context 24{,}102 passed the gate and died at
boot. \textbf{Fix:} the overhead term now uses the resolved launch
concurrency (regression test pins the exact scenario above). At user-set
$m{=}512$/$1024$ the closed gap would have been 1.0/1.8\gib.

\subsection{False-accept: unpinned context certified at an invented 4096}

If neither the request nor the catalog pinned \code{--max-model-len}, the gate
certified context 4096 --- but vLLM boots at the model's \emph{native}
\code{max\_position\_embeddings} (e.g.\ 131k), so the certificate was for a
contract weaker than reality. \textbf{Fix:} an unpinned context now validates
against the native value. Related: user free-form \code{vllm\_args} (which
reach vLLM \emph{last} and win) were invisible to the gate and could raise the
context after certification; the gate now parses them at highest precedence.

\subsection{Availability: $\sim$20 curated VLM models were unlaunchable}

VLM configs nest the language model's dimensions under
\code{text\_config}/\code{llm\_config}/\code{language\_config}; the metadata
mapper read only top-level keys, so every such model resolved to ``cannot
verify'' and was \emph{hard-blocked} --- a regression for models that launch
fine today. \textbf{Fix (two parts):} the mapper now merges the nested
language-model sub-config (verified against real InternVL2.5 and
Llama-3.2-Vision--shaped configs), and verdicts that mean ``we cannot model
this'' are flagged \code{unverifiable} and \emph{skip} the gate instead of
blocking. Models whose configs are irreducibly sparse (llava-1.5's
\code{text\_config} relies on transformers class defaults we refuse to
replicate; deepseek-vl2; glm-4v; Pixtral has no \code{config.json}) launch
unvalidated with a logged warning --- exactly the pre-gate behavior.

\subsection{Availability: multi-node entries were hard-blocked}

Eight curated entries (\code{Llama-405B}, Mistral-Large, Mixtral-8x22B,
Command-R+, Llama-3.2-90B-Vision) have \code{num\_nodes}$>1$ and were
unconditionally rejected because multi-node memory math is out of scope.
\textbf{Fix:} multi-node launches skip the gate with a warning. (Policy note:
the team may later prefer advisory warnings in the UI; the math for pipeline
parallelism remains future work.)

\subsection{Second review round: adversarial findings}

An adversarial pass over the fixed code (trying to \emph{break} the gate)
surfaced and closed a further set of defects:

\begin{itemize}[leftmargin=1.4em]
\item \textbf{Comma-joined flag strings.} LLMHub's wire format joins vLLM
  flags with commas (\code{--max-model-len=131072,--max-num-seqs=32}); the
  flag parser was whitespace-delimited, so only the \emph{last} flag of a
  multi-flag string parsed --- re-opening the user-override false accept for
  every earlier flag. The parser now splits on both.
\item \textbf{Model-identity steering.} The request's \code{hf\_model} field
  steered gate sizing, but vec-inf launches the \emph{catalog} model's local
  weights when they exist --- a tiny stand-in repo could earn a confident
  \code{valid=True} for a 70B launch. Sizing now follows the catalog-derived
  identity whenever the two disagree (logged).
\item \textbf{Unmodeled memory flags.} Free-form flags that change vLLM's
  memory picture (\code{--gpu-memory-utilization}, \code{--dtype},
  \code{--kv-cache-dtype}, LoRA/speculative flags, \dots) now cause a loud
  skip: certifying a config whose memory behavior we did not model would be a
  guess dressed up as a verdict.
\item \textbf{Dequantized-at-load checkpoints.} MXFP4/FP8 checkpoints (e.g.\
  \code{gpt-oss-120b}) have no native kernels on Ampere; vLLM materializes
  bf16 weights 2--4$\times$ the on-disk size we measure. On A40/A100 such
  models are now refused as unverifiable (skip) instead of certified at disk
  size and OOMing at load. On H200 they size normally.
\item \textbf{MLA models.} DeepSeek-V2/V3--family configs (compressed latent
  KV) are over-counted $\sim$10--25$\times$ by the standard formula --- a
  confident wrong verdict in either direction --- and are now refused as
  unverifiable.
\item \textbf{Crash-to-bypass hardening.} A crash inside the gate skips it
  (by design), so crashing must never be the cheap way past: concurrency
  overrides are clamped to $\ge 1$, all catalog/flag integer coercions are
  guarded (a YAML \code{--max-model-len: yes} no longer certifies a context
  of one token), and mixed-provenance VLM configs (top-level vision dims
  combined with partial text dims) can no longer produce a confidently wrong
  \code{head\_dim}.
\item \textbf{Launch-path regression (pre-existing).} \code{num\_nodes} always
  arrives as an explicit \code{None} from the request schema, so the old
  dict-get default never fired and \code{num\_gpus $\times$ num\_nodes} raised
  \code{TypeError} on every dialog launch that set a GPU count. Fixed and
  covered by new service-level tests.
\end{itemize}

The same pass produced strong negative evidence: the resolved-concurrency
threading is correct end to end, the UI's \code{starts} and the gate's verdict
were shown \emph{algebraically identical} and agree on a 900-point sweep plus
a boundary-crossing regression test, TP edge cases (replication floor,
non-divisible heads, partition caps) could not be made optimistic, and the
streamed header reads could not be made to buffer unbounded data.

\subsection{Robustness: crash paths, HF coupling, H200}

The gate ran up to four sequential HTTPS calls to Hugging Face per launch
(after GPU allocation, uncached), an unexpected estimator exception would
surface as a 500 \emph{and leak the GPU-allocation counter}, and a
range-ignoring mirror could stream a multi-GB weights file into backend
memory. \textbf{Fixes:} successful metadata lookups are cached (10-minute
TTL, bounded size; failures are never cached), the gate call is wrapped
fail-open with the exception logged, safetensors header reads are streamed
with hard byte limits, malformed headers degrade to ``unknown weights''
(never a partial undercount, which would be optimistic), and 401/403 from HF
now produce actionable messages (missing token vs.\ unaccepted license).
H200 rows carried the nominal 141\gib{} although H200s report
$\approx$140.4\gib{} in \code{nvidia-smi}; the table now carries a padded
educated guess of \textbf{140.0\gib} until probes run on real hardware.

\section{Caveats and known limitations}

\begin{itemize}[leftmargin=1.4em]
\item \textbf{H200 is an educated guess.} Neither its VRAM nor the
  internal-overhead line has been probed there ($0.4\gib$ pad applied);
  same for TP$\ge$8 (the 0.25\gib{} TP buffer is a safety margin, measured
  $\approx$0 at TP$\le$4). First H200/TP-8 probes should tighten both.
\item \textbf{The Option-B surface is intentional.} The gate certifies boot,
  not saturated concurrency. A booted model under extreme burst degrades by
  queueing/preemption, not OOM; capacity numbers in the UI are advisory.
\item \textbf{Unverifiable models launch unvalidated} (llava-1.5/1.6,
  deepseek-vl2, glm-4v, Pixtral, plus anything during an HF outage). This is
  a deliberate availability-over-strictness call for the \emph{launch gate}
  only; every skip is logged. The team should confirm this policy.
\item \textbf{Workload archetype factors are placeholders} (marked as such in
  the data file); they affect only the advisory ``typical'' column, never
  gating.
\item \textbf{Architectures the formula over-estimates:} MLA models
  (DeepSeek-V2/V3 compress KV; our formula can over-charge them
  $\sim$25$\times$ --- they are multi-node anyway, hence skipped) and
  sliding-window/hybrid-attention models (Mistral, Gemma-2/3) where newer
  vLLM allocates less than full KV per layer. Both err conservative.
  FP8-KV runs are likewise over-charged $2\times$ (we use the weight dtype
  for KV). None of these can cause a false accept.
\item \textbf{SU rate anomaly:} \code{gpuA100x8-interactive} is tabled at
  3000 SU/GPU-hr while every other interactive partition is exactly $2\times$
  its batch rate (which would give 2000). Cost display only --- gating is
  unaffected --- but worth confirming against Delta billing.
\item \textbf{Calibration breadth:} eight probes, two models ($\le$7B), two
  GPU types, one vLLM version (0.11.0). The conservative-line + zero-slack
  regression test protects against silent drift, but two extrapolations are
  load-bearing and unprobed: 70B-class verdicts (vLLM's profile-run peak
  scales with \code{hidden\_size}, which the flat overhead line measured on
  0.5B/7B does not capture --- some passing 70B configs sit within
  $\sim$1\gib{} of the line) and any vLLM upgrade on Delta (the overhead line
  is version-dependent). A 70B/H200/TP-8 probe round is the single highest
  value follow-up.
\item \textbf{The two API surfaces are deliberately different.} The launch
  gate certifies boot; \code{POST /api/validate-config} is a strict
  worst-case survey (KV at $256\times$ context) and will say
  \code{valid:false} for many configs that boot and launch fine. Responses
  now carry \code{unverifiable} so callers can tell ``will not fit'' from
  ``cannot model it (and the launcher would proceed)''. The survey's
  \code{fits} field is the same fictional worst case --- UI and tooling should
  read \code{starts}.
\item \textbf{Operational prerequisites.} \code{HF\_TOKEN} must be set in the
  production backend environment: without it every gated repo
  (\code{meta-llama/*}, \code{google/*}, some \code{mistralai/*} ---
  roughly half the catalog) resolves as unverifiable and launches
  \emph{ungated}. Skips are logged at WARNING; a periodic check of skip rate
  and skip reasons is the cheapest way to notice the gate silently losing
  coverage.
\end{itemize}

\section{Live end-to-end validation on Delta}
\label{sec:live}

The full path was exercised live from a Delta login node against a throwaway
local database (the shared database was never touched):

\begin{itemize}[leftmargin=1.4em]
\item \textbf{Gate-block, live}: a deliberate over-budget request
  (Qwen2.5-7B at context 500{,}000 on one A40) returned HTTP~400
  ``Launch blocked: \dots\ over by 1.7\gib'' with the split-contract
  breakdown, released the GPU allocation, and submitted no Slurm job.
\item \textbf{Launch, live}: Qwen2.5-0.5B-Instruct through
  \code{POST /api/models/deployments} $\to$ gate pass $\to$ vec-inf $\to$
  Slurm job 21145772 $\to$ \emph{vLLM booted} (``Application startup
  complete'') on \code{gpuA40x4}; the job was then cancelled.
\item \textbf{Model validated against the boot log}: weights loaded at
  0.9271\gib{} (fixture: 0.9267); the KV pool was 3{,}288{,}608 tokens at
  37.64\gib{} --- \emph{exactly} 12{,}288 bytes/token, matching the formula
  to the byte.
\item \textbf{One real defect found and fixed}: the job booted with
  \code{max\_num\_seqs}${}=1024$ --- the vLLM \emph{V1-engine} default ---
  while the resolver assumed the V0-era 256. Measured internal overhead was
  1.92\gib{} vs.\ the 1.31\gib{} the gate would have charged at 256
  ($\sim$0.6\gib{} optimistic); at the correct 1024 the model charges
  2.85\gib{} (conservative by 0.93\gib). \code{DEFAULT\_MAX\_NUM\_SEQS}
  is now 1024, and the live measurement is vendored as a ninth calibration
  probe so regression tests pin it. (The frontend's own copy of the 256
  constant needs the same update in the frontend pass.)
\end{itemize}

\section{Verification status}

133 backend tests pass (including the live production-container probe) (fully offline; real HF configs and the eight probe
measurements are vendored as fixtures), including regression tests for the
boot-window false accept, comma-joined flag parsing, model-identity
protection, the skip semantics, nested-config resolution and its
mixed-provenance guard, MLA/dequantization refusals, metadata-cache behavior
(TTL, no negative caching, bounded eviction), malformed-safetensors handling,
survey-vs-gate agreement across the boot boundary, and the
\code{launch\_model} resource-accounting contract (release on block, fail-open
on crash, the \code{num\_nodes} regression). The gate-pass/boot-fail scenario,
user-override visibility, native-context fallback, and both skip paths were
additionally exercised end-to-end against the real gate code with only the
network fetch stubbed, and two independent adversarial review passes were run
over the final code.

\end{document}
